\documentclass[11pt,a4paper]{article}
\usepackage[margin=1in]{geometry}
\usepackage[hidelinks]{hyperref}
\usepackage{longtable}
\usepackage{booktabs}
\usepackage{array}
\usepackage{listings}
\usepackage{xcolor}
\usepackage{xurl}
\usepackage[T1]{fontenc}
\usepackage[utf8]{inputenc}

\lstset{
  basicstyle=\ttfamily\small,
  breaklines=true,
  breakatwhitespace=false,
  breakindent=0pt,
  frame=single,
  columns=fullflexible,
  xleftmargin=0pt,
  framexleftmargin=6pt,
  framexrightmargin=6pt,
  aboveskip=6pt,
  belowskip=6pt
}
\setlength{\emergencystretch}{3em}
\Urlmuskip=0mu plus 2mu

\title{newdb Shared Bundle External Integration Manual (Pure LaTeX)}
\author{newdb Project}
\date{\today}

\begin{document}
\maketitle
\tableofcontents
\newpage

\section{Purpose and Scope}
This document is a pure LaTeX technical manual for integrating \texttt{newdb}
shared libraries from external language runtimes. It covers:
\begin{itemize}
  \item Build and distribution model for \texttt{shared\_bundle}
  \item Full C API and gtest C API references
  \item Command-driven execution model and templates
  \item Unified result parsing and error handling
  \item Deployment and operational recommendations
\end{itemize}

\section{Shared Bundle Model}
\subsection{Output directories}
\begin{itemize}
  \item VS/MSVC: \texttt{build/newdb-vs2026-mt-tests/shared\_bundle/Release}
  \item MinGW: \texttt{build/newdb-mingw/shared\_bundle/Release}
  \item WSL/Linux: \texttt{build/newdb-wsl-tests/shared\_bundle/Release}
\end{itemize}

\subsection{Build commands}
\begin{lstlisting}
# VS/MSVC
cmake --build build/newdb-vs2026-mt-tests --config Release --target shared_bundle --parallel

# MinGW
cmake --build build/newdb-mingw --target shared_bundle --parallel

# WSL/Linux
cmake --build build/newdb-wsl-tests --target shared_bundle -j8
\end{lstlisting}

\subsection{Dependency precedence policy}
\begin{itemize}
  \item Windows bundles include \texttt{run\_*.cmd} launchers to prioritize local DLLs.
  \item Linux bundles are linked with \texttt{RPATH=\$ORIGIN} to prefer colocated \texttt{.so}.
  \item MinGW bundle includes runtime DLLs:
  \begin{itemize}
    \item \texttt{libstdc++-6.dll}
    \item \texttt{libgcc\_s\_*.dll}
    \item \texttt{libwinpthread-1.dll}
  \end{itemize}
\end{itemize}

\section{newdb C API Reference}
\subsection{Header and compatibility constants}
Header path:
\begin{itemize}
  \item \texttt{newdb/engine/include/newdb/c\_api.h}
\end{itemize}

Version/ABI constants:
\begin{itemize}
  \item \texttt{NEWDB\_C\_API\_VERSION\_MAJOR/MINOR/PATCH}
  \item \texttt{NEWDB\_C\_API\_ABI\_VERSION}
\end{itemize}

\subsection{Error codes}
\begin{longtable}{>{\ttfamily}p{0.38\linewidth}p{0.54\linewidth}}
\toprule
Code & Description \\
\midrule
NEWDB\_OK & Success \\
NEWDB\_ERR\_INVALID\_ARGUMENT & Invalid argument or buffer contract violation \\
NEWDB\_ERR\_INVALID\_HANDLE & Invalid or destroyed session handle \\
NEWDB\_ERR\_EXECUTION\_FAILED & Command/business execution failed \\
NEWDB\_ERR\_INTERNAL & Internal engine failure \\
NEWDB\_ERR\_LOG\_IO & Log or file I/O failure \\
NEWDB\_ERR\_SESSION\_TERMINATED & Session terminated and cannot continue \\
\bottomrule
\end{longtable}

\subsection{Thread safety contract}
\begin{itemize}
  \item Global metadata APIs are thread-safe.
  \item Per-session APIs are safe only when one thread owns one handle.
  \item Concurrent use of the same handle is not guaranteed safe.
\end{itemize}

\subsection{Exported functions}
\subsubsection*{Global/metadata}
\begin{itemize}
  \item \texttt{newdb\_version\_string}
  \item \texttt{newdb\_api\_version\_major}, \texttt{newdb\_api\_version\_minor},
  \texttt{newdb\_api\_version\_patch}
  \item \texttt{newdb\_abi\_version}
  \item \texttt{newdb\_negotiate\_abi}
  \item \texttt{newdb\_error\_code\_string}
  \item \texttt{newdb\_sum}
\end{itemize}

\subsubsection*{Schema/session/execution}
\begin{itemize}
  \item \texttt{newdb\_check\_schema\_file}
  \item \texttt{newdb\_session\_create}
  \item \texttt{newdb\_session\_destroy}
  \item \texttt{newdb\_session\_set\_table}
  \item \texttt{newdb\_session\_last\_error}
  \item \texttt{newdb\_session\_execute}
  \item \texttt{newdb\_session\_runtime\_stats}
  \item \texttt{newdb\_session\_append\_runtime\_snapshot}
\end{itemize}

\subsection{Recommended lifecycle}
\begin{enumerate}
  \item Load library by absolute path.
  \item Verify ABI via \texttt{newdb\_negotiate\_abi}.
  \item Create session.
  \item Execute command stream.
  \item Capture runtime stats/snapshots if needed.
  \item Destroy session.
\end{enumerate}

\section{gtest\_capi Reference}
Header path:
\begin{itemize}
  \item \texttt{newdb/include/gtest\_capi.h}
\end{itemize}

Main API groups:
\begin{itemize}
  \item Init/Run:
  \texttt{gtest\_capi\_init\_from\_argv},
  \texttt{gtest\_capi\_run\_all}
  \item Discovery:
  \texttt{gtest\_capi\_list\_tests},
  \texttt{gtest\_capi\_enumerate\_tests}
  \item Options:
  filter/output/repeat/shuffle/random\_seed/color and failure/exception controls
  \item Generic option bridge:
  \texttt{gtest\_capi\_set\_option},
  \texttt{gtest\_capi\_get\_option}
  \item Counters:
  suite/test totals, pass/fail/skip/disabled and elapsed time
\end{itemize}

\section{Feature-Oriented Integration Guidance}
Through \texttt{newdb\_session\_execute}, external systems can drive:
\begin{itemize}
  \item DDL/schema management
  \item DML and query operations
  \item Transaction and conflict/isolation controls
  \item WAL/recovery controls
  \item Runtime maintenance controls (vacuum, group commit, hot index)
\end{itemize}

Recommended wrapper layers:
\begin{itemize}
  \item DDL wrapper
  \item DML wrapper
  \item Query wrapper
  \item Txn wrapper
  \item Runtime tuning wrapper
\end{itemize}

\section{Command Templates for \texttt{newdb\_session\_execute}}
\subsection{DDL and schema}
\begin{lstlisting}
CREATE TABLE(users)
USE(users)
SET PRIMARY KEY(uid)
CREATE SCHEMA(finance_v1)
DROP SCHEMA(finance_v1)
ALTER TABLE users SET SCHEMA(finance_v1)
ALTER TABLE users REMOVE SCHEMA
RENAME TABLE(users_v2)
DROP TABLE(users_v2)
\end{lstlisting}

\subsection{DML and query}
\begin{lstlisting}
DEFATTR(name:string, balance:int, dept:string)
INSERT(1, alice, 1000, rd)
UPDATE(1, alice, 1200)
DELETE(1)
DELETEPK(alice)
SETATTR(1, dept, infra)
DELATTR(dept)
RENATTR(balance, amount)
FIND(1)
WHERE(balance, >, 1000, AND, dept, =, rd)
COUNT
PAGE(1, 20, id, desc)
SUM(balance)
AVG(balance)
MIN(balance)
MAX(balance)
\end{lstlisting}

\subsection{Transaction and runtime tuning}
\begin{lstlisting}
BEGIN
SAVEPOINT(sp1)
ROLLBACK TO(sp1)
RELEASE SAVEPOINT(sp1)
COMMIT
TXNISOLATION snapshot
TXNISOLATION read_committed
WRITECONFLICT reject
WRITECONFLICT wait 2000
AUTOVACUUM threshold 500
AUTOVACUUM interval 60
WALSYNC normal 200
WALADAPTIVE on
GROUPCOMMIT 10 128
HOTINDEX on
\end{lstlisting}

\section{Unified Result Parsing Standard}
Always evaluate:
\begin{enumerate}
  \item Return code (\texttt{rc}) from \texttt{newdb\_session\_execute}
  \item \texttt{newdb\_session\_last\_error}
  \item Human-readable \texttt{output\_buf}
\end{enumerate}

Recommended result DTO:
\begin{lstlisting}
ExecuteResult {
  rc: int
  last_error: int
  error_name: string
  output: string
  ok: bool
}
\end{lstlisting}

Buffer sizing guideline:
\begin{itemize}
  \item 4KB minimum for simple commands
  \item 16KB+ for complex query output
\end{itemize}

\section{Deployment and Operations}
\subsection{Deployment checklist}
\begin{enumerate}
  \item Distribute full \texttt{shared\_bundle/Release} directory.
  \item Use absolute-path dynamic loading.
  \item Verify ABI before first session creation.
  \item One thread per session handle.
  \item Record \texttt{rc + last\_error + output} for all commands.
\end{enumerate}

\subsection{Observability}
Collect:
\begin{itemize}
  \item Commit latency and conflict metrics
  \item WAL recovery elapsed/scanned records
  \item Vacuum trigger count and reclaimed bytes
\end{itemize}

\section{Troubleshooting}
\begin{itemize}
  \item ``Module not found'': missing colocated dependencies in bundle.
  \item Wrong DLL/SO loaded: enforce absolute path and launcher usage.
  \item WSL googletest download issues: use local
  \texttt{FETCHCONTENT\_SOURCE\_DIR\_GOOGLETEST}.
\end{itemize}

\end{document}

